% GENERATED FILE. Edit canonical lessons, then run npm run generate:books.
% canonical-source-hash: fnv1a64:742f7c8bfe2173c9
% canonical-lessons: KA-C514-pathya, KA-C514-jivasatva, KA-C514-soragu, KA-C514-dani, KA-C514-ayusya

\chapter{Diet and Strength}
\label{ch:ka-diet-and-strength}
\hlchaptermodality{514}

\begin{chapteropening}
\textbf{By the end of this chapter:} \emph{I can say a diet, a vitamin, to waste away, to get tired and lifespan.}

Complete the last lesson of chapter 514: I can say a diet, a vitamin, to waste away, to get tired and lifespan.
\end{chapteropening}

\section[\texorpdfstring{\kn{ಪಥ್ಯ} (pathya)}{ಪಥ್ಯ (pathya)}]{\kn{ಪಥ್ಯ} (pathya) — a diet; a special diet}

\label{lesson:KA-C514-pathya}



\begin{quote}
Before the new one: say the Kannada for nursing, then the Kannada for a herbal decoction.
\end{quote}

\subsection*{You'll want to know: \kn{ಪಥ್ಯ}}
\textbf{\kn{ಪಥ್ಯ}} — \emph{pathya} — \textquotedblleft{}a diet; a special diet\textquotedblright{}.

\begin{quote}
\kn{ಮಧುಮೇಹ} \kn{ಇರುವುದರಿಂದ} \kn{ಅಪ್ಪ} \kn{ಪಥ್ಯ} \kn{ಮಾಡುತ್ತಾರೆ}.
\end{quote}

\begin{quote}
\emph{madhumēha iruvudarinda appa pathya māḍuttāre.}
\end{quote}

\begin{quote}
Because he has diabetes, father keeps to a diet.
\end{quote}

\kn{ಪಥ್ಯ} \kn{ಮಾಡು}, to keep to a diet.

\begin{cousinweb}
\textbf{\kn{ಪಥ್ಯ}} (\emph{pathya}) is borrowed from Sanskrit \emph{pathya}, \textquotedblleft{}suitable\textquotedblright{} (CDIAL 7746): food that suits the patient.
\end{cousinweb}

\begin{grammarlens}[title={What you've built}]
One more word for describing what you lived and felt.
\end{grammarlens}

\subsection*{Your turn}
\begin{itemize}
\raggedright
  \item \emph{Say it:} \emph{pathya}
  \item \emph{Say it:} \emph{pathya}, once more
  \item \emph{Say it:} \emph{madhumēha iruvudarinda appa pathya māḍuttāre.} — Because he has diabetes, father keeps to a diet.
  \item \emph{Recall:} say the Kannada for nursing, then the Kannada for a herbal decoction, then say \emph{pathya} again
\end{itemize}

\begin{tcolorbox}[breakable,colback=teal!4,colframe=teal!35!black,title={Before you move on}]
What does \kn{ಪಥ್ಯ} mean? (\textquotedblleft{}A diet; a special diet\textquotedblright{}.) Say it once more.
\end{tcolorbox}
\section[\texorpdfstring{\kn{ಜೀವಸತ್ವ} (jīvasatva)}{ಜೀವಸತ್ವ (jīvasatva)}]{\kn{ಜೀವಸತ್ವ} (jīvasatva) — a vitamin}

\label{lesson:KA-C514-jivasatva}



\begin{quote}
Before the new one: say the Kannada for a herbal decoction, then the Kannada for a diet.
\end{quote}

\subsection*{You'll want to know: \kn{ಜೀವಸತ್ವ}}
\textbf{\kn{ಜೀವಸತ್ವ}} — \emph{jīvasatva} — \textquotedblleft{}a vitamin\textquotedblright{}.

\begin{quote}
\kn{ಹಣ್ಣುಗಳಲ್ಲಿ} \kn{ತುಂಬಾ} \kn{ಜೀವಸತ್ವಗಳು} \kn{ಇವೆ}.
\end{quote}

\begin{quote}
\emph{haṇṇugaḷalli tumbā jīvasatvagaḷu ive.}
\end{quote}

\begin{quote}
Fruits have lots of vitamins.
\end{quote}

Plural: \kn{ಜೀವಸತ್ವಗಳು}, \emph{jīvasatvagaḷu}.

\kn{ಜೀವ}, life + \kn{ಸತ್ವ}, essence. Everyday talk also says \kn{ವಿಟಮಿನ್}.

\begin{cousinweb}
\textbf{\kn{ಜೀವಸತ್ವ}} (\emph{jīvasatva}) is a modern learned compound built from Sanskrit \emph{jīva}, \textquotedblleft{}life\textquotedblright{} (CDIAL 5239), and \emph{sattva}, \textquotedblleft{}being, true character\textquotedblright{} (CDIAL 13111).
\end{cousinweb}

\begin{grammarlens}[title={What you've built}]
One more word for describing what you lived and felt.
\end{grammarlens}

\subsection*{Your turn}
\begin{itemize}
\raggedright
  \item \emph{Say it:} \emph{jīvasatva}
  \item \emph{Say it:} \emph{jīvasatva}, once more
  \item \emph{Say it:} \emph{haṇṇugaḷalli tumbā jīvasatvagaḷu ive.} — Fruits have lots of vitamins.
  \item \emph{Recall:} say the Kannada for a herbal decoction, then the Kannada for a diet, then say \emph{jīvasatva} again
\end{itemize}

\begin{tcolorbox}[breakable,colback=teal!4,colframe=teal!35!black,title={Before you move on}]
What does \kn{ಜೀವಸತ್ವ} mean? (\textquotedblleft{}A vitamin\textquotedblright{}.) Say it once more.
\end{tcolorbox}
\section[\texorpdfstring{\kn{ಸೊರಗು} (soragu)}{ಸೊರಗು (soragu)}]{\kn{ಸೊರಗು} (soragu) — to waste away; to grow thin}

\label{lesson:KA-C514-soragu}



\begin{quote}
Before the new one: say the Kannada for a diet, then the Kannada for a vitamin.
\end{quote}

\subsection*{You'll want to know: \kn{ಸೊರಗು}}
\textbf{\kn{ಸೊರಗು}} — \emph{soragu} — \textquotedblleft{}to waste away; to grow thin\textquotedblright{}.

\begin{quote}
\kn{ಕಾಯಿಲೆಯಿಂದ} \kn{ಅವನು} \kn{ತುಂಬಾ} \kn{ಸೊರಗಿದನು}.
\end{quote}

\begin{quote}
\emph{kāyileyinda avanu tumbā soragidanu.}
\end{quote}

\begin{quote}
He grew very thin from the illness.
\end{quote}

Past (I …): \kn{ಸೊರಗಿದೆ}, \emph{soragide}.

Of people, plants and animals.

\begin{cousinweb}
\textbf{\kn{ಸೊರಗು}} (\emph{soragu}) is native Dravidian: DEDR 2687 lists Kannada \emph{soragu}, \textquotedblleft{}to wither, grow dry, grow lean\textquotedblright{}, beside Tamil \emph{curuṅku}, \textquotedblleft{}to shrink\textquotedblright{}.
\end{cousinweb}

\begin{grammarlens}[title={What you've built}]
One more word for describing what you lived and felt.
\end{grammarlens}

\subsection*{Your turn}
\begin{itemize}
\raggedright
  \item \emph{Say it:} \emph{soragu}
  \item \emph{Say it:} \emph{soragu}, once more
  \item \emph{Say it:} \emph{kāyileyinda avanu tumbā soragidanu.} — He grew very thin from the illness.
  \item \emph{Recall:} say the Kannada for a diet, then the Kannada for a vitamin, then say \emph{soragu} again
\end{itemize}

\begin{tcolorbox}[breakable,colback=teal!4,colframe=teal!35!black,title={Before you move on}]
What does \kn{ಸೊರಗು} mean? (\textquotedblleft{}To waste away; to grow thin\textquotedblright{}.) Say it once more.
\end{tcolorbox}
\section[\texorpdfstring{\kn{ದಣಿ} (daṇi)}{ದಣಿ (daṇi)}]{\kn{ದಣಿ} (daṇi) — to get tired}

\label{lesson:KA-C514-dani}



\begin{quote}
Before the new one: say the Kannada for a vitamin, then the Kannada for to waste away.
\end{quote}

\subsection*{You'll want to know: \kn{ದಣಿ}}
\textbf{\kn{ದಣಿ}} — \emph{daṇi} — \textquotedblleft{}to get tired\textquotedblright{}.

\begin{quote}
\kn{ಇಡೀ} \kn{ದಿನ} \kn{ಕೆಲಸ} \kn{ಮಾಡಿ} \kn{ಅವನು} \kn{ದಣಿದನು}.
\end{quote}

\begin{quote}
\emph{iḍī dina kelasa māḍi avanu daṇidanu.}
\end{quote}

\begin{quote}
He got tired after working all day.
\end{quote}

Past (I …): \kn{ದಣಿದೆ}, \emph{daṇide}.

\kn{ದಣಿವು}, tiredness; \kn{ದಣಿವಾರಿಸಿಕೊ}, to rest a while.

\begin{cousinweb}
\textbf{\kn{ದಣಿ}} (\emph{daṇi}) is native Dravidian: DEDR 3045 lists Kannada \emph{daṇi}, \emph{daṇivu}, beside \emph{taṇi}, \textquotedblleft{}to grow cool, get feeble or fatigued\textquotedblright{}.
\end{cousinweb}

\begin{grammarlens}[title={What you've built}]
One more word for describing what you lived and felt.
\end{grammarlens}

\subsection*{Your turn}
\begin{itemize}
\raggedright
  \item \emph{Say it:} \emph{daṇi}
  \item \emph{Say it:} \emph{daṇi}, once more
  \item \emph{Say it:} \emph{iḍī dina kelasa māḍi avanu daṇidanu.} — He got tired after working all day.
  \item \emph{Recall:} say the Kannada for a vitamin, then the Kannada for to waste away, then say \emph{daṇi} again
\end{itemize}

\begin{tcolorbox}[breakable,colback=teal!4,colframe=teal!35!black,title={Before you move on}]
What does \kn{ದಣಿ} mean? (\textquotedblleft{}To get tired\textquotedblright{}.) Say it once more.
\end{tcolorbox}
\section[\texorpdfstring{\kn{ಆಯುಷ್ಯ} (āyuṣya)}{ಆಯುಷ್ಯ (āyuṣya)}]{\kn{ಆಯುಷ್ಯ} (āyuṣya) — lifespan; long life}

\label{lesson:KA-C514-ayusya}



\begin{quote}
Before the new one: say the Kannada for to waste away, then the Kannada for to get tired.
\end{quote}

\subsection*{You'll want to know: \kn{ಆಯುಷ್ಯ}}
\textbf{\kn{ಆಯುಷ್ಯ}} — \emph{āyuṣya} — \textquotedblleft{}lifespan; long life\textquotedblright{}.

\begin{quote}
\kn{ಒಳ್ಳೆಯ} \kn{ಆಹಾರದಿಂದ} \kn{ಆಯುಷ್ಯ} \kn{ಹೆಚ್ಚುತ್ತದೆ}.
\end{quote}

\begin{quote}
\emph{oḷḷeya āhāradinda āyuṣya heccuttade.}
\end{quote}

\begin{quote}
Good food lengthens life.
\end{quote}

Elders bless with \kn{ನೂರು} \kn{ವರ್ಷ} \kn{ಆಯುಷ್ಯ}, a hundred years of life.

\begin{cousinweb}
\textbf{\kn{ಆಯುಷ್ಯ}} (\emph{āyuṣya}) is borrowed from Sanskrit \emph{āyuṣya}, \textquotedblleft{}giving long life; long life\textquotedblright{} (CDIAL 1292a).
\end{cousinweb}

\begin{grammarlens}[title={What you've built}]
One more word for describing what you lived and felt.
\end{grammarlens}

\subsection*{Your turn}
\begin{itemize}
\raggedright
  \item \emph{Say it:} \emph{āyuṣya}
  \item \emph{Say it:} \emph{āyuṣya}, once more
  \item \emph{Say it:} \emph{oḷḷeya āhāradinda āyuṣya heccuttade.} — Good food lengthens life.
  \item \emph{Recall:} say the Kannada for to waste away, then the Kannada for to get tired, then say \emph{āyuṣya} again
\end{itemize}

\begin{tcolorbox}[breakable,colback=teal!4,colframe=teal!35!black,title={Before you move on}]
What does \kn{ಆಯುಷ್ಯ} mean? (\textquotedblleft{}Lifespan; long life\textquotedblright{}.) Say it once more.
\end{tcolorbox}
